\documentclass[../../../include/open-logic-section]{subfiles}

\begin{document}

\olfileid{sth}{z}{arbintersections}
\olsection{ضمیمہ۔ بندش، خاصیت توں سیٹ بنانا تے اشتراک}

\olref[sth][z][milestone]{sec} وچ اسیں آکھیا سی کہ تسی
\olref[sfr][][]{part} دے سادے طریقے نال کیتے کم ول مُڑ کے ویکھو تے
پرکھو کہ اوہ~$\Zminus$ وچ کیتا جا سکدا اے۔ جے تسی ایویں کیتا، تاں
\emph{اک} گل وچ اٹک سکدے سی۔ ڈیڈیکنڈ دے \emph{بندشاں} والے بیان وچ
\emph{اشتراک} دی ورتوں۔

\olref[sfr][infinite][dedekind]{Closure} توں یاد کرو کہ
\begin{align*}
  \closureofunder{f}{o} & = \bigcap\Setabs{X}{o \in X \text{ تے $X$، $f$ تحت
بند اے}}.
\intertext{ایہدی عام صورت ایہو جیہی تعریف اے:}
  C & = \bigcap\Setabs{X}{\phi(X)}.
\end{align*}
پر ایس گل نال سانوں خبردار ہو جانا چاہیدا اے۔ سادہ خاصیت توں سیٹ بناؤن
دا اصول نہیں چلدا، سو کوئی ضمانت نہیں کہ $\Setabs{X}{\phi(X)}$ موجود اے۔
فیر ایویں لگدا اے کہ ایہو جیہیاں تعریفاں وچ شاید اسیں \emph{ہیرا پھیری}
کر رہے آں۔

چنگی گل ایہ اے کہ ایہ ہیرا پھیری نہیں۔ یا، جے ایہ ایس صورت وچ
\emph{واقعی} ہیرا پھیری اے، تاں اسیں تھوڑی جیہی سچّی محنت نال ایہناں نوں ٹھیک کر
سکدے آں۔ اوس محنت دی جھلک \olref[sth][z][sep]{prop:intersectionsexist}
وچ سی، جدوں اسیں سمجھایا سی کہ $\bigcap A$ کسے وی $A \neq \emptyset$
لئی کیوں موجود اے۔ پر ہُن اسیں ایہ گل کھول کے دساں گے۔

عنصراں نال تعیّن دے اصول نال، جے اسیں $C$ نوں
$\bigcap\Setabs{X}{\phi(X)}$ دے طور اُتے تعریف کرن دی کوشش کریے،
تاں اصل وچ سانوں بس اک ایسی شے $C$ چاہیدی اے جیہڑی ایہ شرط پوری کرے:
\begin{equation}\ollabel{bicondelimarbintersection}
	\forall x(x \in C \liff \forall X(\phi(X) \lif x \in X))\tag{*}
\end{equation}
ہُن منّ لو \emph{کوئی} سیٹ $S$ اے جیہدے لئی $\phi(S)$ اے۔ فیر
\olref{bicondelimarbintersection} پوری کرن لئی اسیں $C$ نوں
\emph{وکھ کرن دے اصول} نال ایویں تعریف کر سکدے آں:
\[
	C = \Setabs{x \in S}{\forall X(\phi(X) \lif x \in X)}.
\]
ایہ پرکھنا مشق چھڈدے آں کہ ایس تعریف نال \olref{bicondelimarbintersection}
ملدی اے، جیویں سانوں چاہیدی سی۔
%  fix $x$. First suppose that $x \in C$, as (re)defined; then both $x
%  \in S$ and $\forall X(\phi(X) \lif x \in X)$; but since
%  $\phi(S)$, this reduces to the condition that $\forall X(\phi(X)
%  \lif x \in X)$. Conversely, suppose $\forall X(\phi(X) \lif
%  x \in X)$; then, since $\phi(S)$, we also have that $x \in S$; so
%  $x \in c$, as (re)defined. 
ایہ عام طریقہ سانوں اشتراکاں دی تعریف وچ سادہ خاصیت توں سیٹ بناؤن دے
اصول دی ہر ظاہری ورتوں توں بچا دیوے گا۔ جیہڑی خاص صورت نال ایہ گل
شروع ہوئی، یعنی $\closureofunder{f}{o}$، اوہدے لئی ایہ طریقہ ایویں چلے گا۔
اسیں \olref[sfr][infinite][dedekind]{closureproperties} دا ثبوت ایہ گل
ویکھ کے شروع کیتا سی کہ $o \in \ran{f}\cup\{o\}$ تے $\ran{f} \cup \{o\}$،
$f$ تحت بند اے۔ سو جیہڑی شے سانوں چاہیدی اے، اوہدی تعریف ایویں کر سکدے آں:
\[
	\closureofunder{f}{o} = \Setabs{x \in \ran{f} \cup \{o\}}{(\forall X \ni o)(X \text{ $f$ تحت بند اے} \lif x \in X)}.
\]

\end{document}
